\section{Summary}
\label{sec:summary}

This work defines a flow-based MCMC algorithm to sample lattice field configurations from a desired probability distribution:
\begin{enumerate}
    \item A real NVP flow model is trained to produce approximately the desired distribution;
    \item Samples are proposed from the trained model;
    \item Starting from an arbitrary configuration, each proposal is accepted or rejected to advance a Markov chain using the Metropolis-Hastings algorithm.
\end{enumerate}
The approach is shown to define an ergodic and balanced Markov chain, thus guaranteeing convergence to the desired probability distribution in the limit of a long Markov chain. In essence, the flow-based MCMC algorithm combines the expressiveness of normalizing flows based on neural networks with the theoretical guarantees of Markov chains to create a trainable and asymptotically-correct sampler. Since these flows are applicable for arbitrary configurations with continuous, real-valued degrees of freedom, one can generically apply this method to any of a broad class of lattice theories. Here, the algorithm is implemented in practice for $\phi^4$ theory, and is demonstrated to produce ensembles of configurations that are indistinguishable from those generated using standard local Metropolis and HMC algorithms, based on studies of a number of physical observables.

A key feature of the approach is that models trained to a fixed acceptance rate do not experience critical slowing down in the sampling stage. In particular, the autocorrelation time for all observables is dictated entirely by the accuracy with which the flow model has been trained; perfect training corresponds to decorrelated samples and 100\% acceptance in the Metropolis-Hastings step of the MCMC process. Nevertheless, the efficiency with which the training step of this approach can be scaled to larger model sizes, and to more complicated theories such as QCD, remains to be studied. Recent advances in the training and scaling of flow models provide reasons for optimism on this front. Further, incorporating symmetries generally improves data efficiency of training, and implementing spacetime and gauge symmetries~\cite{Cohen2019GaugeEquiv} may be a natural next step to practically train these flow models for lattice gauge theories like QCD.

In moving towards lattice gauge theories such as QCD, several theoretical developments are also required. The real NVP model chosen to parameterize the normalizing flows here is described in terms of vectors of variables $\phi \in \mathbb{R}^D$. Gauge configurations, however, live in a compact manifold arising from the Lie group structure. Extending this method will require a normalizing flow model that can act on this manifold while remaining sufficiently expressive. The choice of prior likewise will need to be extended to a distribution over the manifold of lattice gauge configurations which can be easily sampled. A uniform distribution, for example, may be a candidate for a prior, but this choice must be tested in the context of a specific flow model.

If the flow-based MCMC algorithm proposed here can be implemented for a complex theory such as QCD, the advantages would be significant; arbitrarily large ensembles of field configurations could be generated at minimal cost. The independence of the proposal step from any previous configuration allows parallel generation of proposals, and the continually-improving support in hardware and software for neural network execution suggests future practical gains for this style of ensemble-generation. Given efficient sample generation from a trained model, ensembles would not need to be stored long-term. Moreover, a model trained for one action could either be re-trained or used as a prior for another flow model targeting an action with nearby parameter values. This would allow efficient tuning of parameters and generation of additional ensembles interpolating between and extrapolating from existing models.
